\chapter{Studio del grafico di una funzione}
Il calcolo differenziale sviluppato fin ora ci permette di affrontare in modo completo un tipico problema, quello di tracciare il grafico di una funzione f.

A tal proposito illustreremo i passaggi fondametali da svolgere:
\begin{enumerate}
    \item Determinare l'insieme su cui f è definita, ossia il dominio di f;
    \item Limiti agli estremi del dominio. Calcolare i limiti destri e sinistri alla frontiera di tale insieme. In questo modo possiamo determinare gli eventuali asintoti orizzontali e verticali e alcuni eventuali punti di discontinuità;
    \item Studio del segno ponendo $f(x)>0$, per determinare gli intervalli in cui la funzione è positiva e quelli in cui è negativa;
    \item Studio della derivata prima per calcolare i massimi e minimi della funzione e la sua crescenza o decrescenza;
    \item Studio della derivata seconda per calcolare le convessità e i punti di flesso;
    \item Disegno del grafico della funzione.
\end{enumerate}